\chapter{The matter ledger of the horizon: the dual-sector model and its chains}\label{ch:dual}

\section{One term added to the horizon entropy}

In the Jacobson and Cai--Kim derivations~\cite{Jacobson1995,CaiKim2005}, the apparent-horizon entropy is geometric, $S_{\mathrm{geo}}=\pi c^3/(G\hbar H^2)$.
IAM adds the accumulated Landauer cost of every irreversible record written on a timelike worldline:
\begin{equation}
 -dE=T_H\,d\!\left(S_{\mathrm{geo}}+S_{\mathrm{info}}\right).
\end{equation}
The information density is taken to grow as $\dot\rho_{\mathrm{info}}=\rho_{\mathrm{info}}H/a$ \conjecture. With $\dot a=aH$ this is
$d\rho_{\mathrm{info}}/da=\rho_{\mathrm{info}}/a^2$, whose solution, normalised to one today, is the activation function \derived:
\begin{equation}
 \boxed{\;E(a)=\exp\!\left(1-\frac1a\right)=e^{-z}\;}
 \label{part2:eq:Ea_chains}
\end{equation}
$E\to0$ at early times, $E(1)=1$, and $E\to e$ as $a\to\infty$. Two properties fix where the effect lives in time:
the inflection of $E(a)$ is at $a=1/2$ ($z=1$), and the production per e-fold, $dE/d\ln a=E/a$, peaks exactly at $a=1$ \calc.
Entropy production per e-fold is at its maximum today.

\section{Two sectors}

Records are written only on timelike worldlines. A null worldline accumulates no proper time ($d\tau=0$), so it writes none.
Matter feels the information term; light does not. In the standard modified-growth form \derived:
\begin{equation}
 \mu(a)=\frac{H^2_{\Lambda\mathrm{CDM}}(a)}{H^2_{\Lambda\mathrm{CDM}}(a)+\beta_m E(a)H_0^2},\qquad \Sigma(a)=1,
 \label{part2:eq:mu}
\end{equation}
where $\mu$ multiplies Newton's constant in the growth equation
$\ddot\delta_m+2H\dot\delta_m-\tfrac32\Omega_mH^2\mu\,\delta_m=0$ and $\Sigma$ multiplies the lensing potential.

\begin{keybox}[title={The coupling is predicted, not fitted}]
The virial theorem gives the matter coupling $\beta_m=\Omega_m/2$ (Chapter~\ref{ch:virial}). With Planck's $\Omega_m=0.3153$~\cite{Planck2018VI},
$\beta_m=0.15765$ \derived. This value is \emph{fixed} in every chain below; nothing in $\beta_m$ is sampled.
An earlier independent check, with $\beta$ left free against expansion-rate and growth data, returned $0.157\pm0.029$ \observed.
\end{keybox}

\begin{checkbox}
\calc\ From Eq.~\eqref{part2:eq:mu} with $\beta_m=0.15765$: $\mu=0.864$ at $z=0$, $0.905$ at $z=0.2$, $0.982$ at $z=1$, $0.998$ at $z=2$.
The coupling $1-\mu_0=13.6\%$ is the change in the effective Newton constant for matter today. It is \emph{not} the change in growth:
growth integrates $\mu$ over time, so the growth factor changes by about $0.8\%$ and $\sigma_8$ by about $1.6\%$ (next section).
\end{checkbox}

\section{The chain record}\label{sec:chains}

Eighteen Markov chains were run against the Planck 2018 likelihood~\cite{Planck2018V,MahaffeyLateTime,MahaffeyDSPerturbation}. All files are in the repository
(\texttt{mgcamb\_validation/\allowbreak{}chains}, \texttt{camb\_validation/\allowbreak{}chains}). The numbers below come from one extraction of the final files
(30\,\% burn-in, weighted posterior means; table \texttt{CHAIN\_EXTRACTION\_FINAL.csv}). Every chain's last recorded Gelman--Rubin~\cite{Gelman1992}
$R-1$ is at or below $0.010$.

\begin{itemize}
 \item \textbf{Level 1, 12 chains (MGCAMB~\cite{Zhao2009MGCAMB,Wang2023MGCAMB}).} Four data sets (Planck; +RSD; +BAO; +Pantheon+), each run three ways: IAM with
       $\beta_m$ fixed at the prediction, $\Lambda$CDM, and a free amplitude $\mu_0$.
 \item \textbf{Level 2, 3 chains (CAMB~\cite{Lewis2000CAMB} with the mechanism written into the perturbation equations).} IAM (Planck), $\Lambda$CDM (Planck), IAM (Planck+RSD).
 \item \textbf{Level 2b, 2 chains.} The same term placed in the background expansion instead of the perturbations.
 \item \textbf{The baryon chain.} Planck with the standard baryon density left free (Chapter~\ref{ch:baryon}).
\end{itemize}

\begin{center}
\small
\begin{tabularx}{\linewidth}{@{}l>{\centering\arraybackslash}X>{\centering\arraybackslash}X>{\centering\arraybackslash}X>{\centering\arraybackslash}X@{}}
\toprule
data & $\Delta\chi^2_{\min}$ (IAM $-$ $\Lambda$CDM) & $\sigma_8$ ($\Lambda$CDM $\to$ IAM) & $H_0$ (IAM) & free $\mu_0$: median (90\,\% lower bound) \\
\midrule
Planck            & $+0.96$ & $0.814\to0.802$ & $67.08\pm0.53$ & $+0.059$ ($-0.305$) \\
Planck + RSD      & $+0.56$ & $0.813\to0.800$ & $67.49\pm0.42$ & $+0.064$ ($-0.204$) \\
Planck + BAO      & $+1.73$ & $0.812\to0.799$ & $67.54\pm0.43$ & $+0.047$ ($-0.304$) \\
Planck + Pantheon+ & $+1.58$ & $0.813\to0.800$ & $67.00\pm0.50$ & $+0.030$ ($-0.342$) \\
\midrule
Level 2, Planck   & $+0.54$ & $0.809\to0.800$ & $67.16\pm0.47$ & -- \\
\bottomrule
\end{tabularx}
\end{center}
$\Delta\chi^2_{\min}$ is the difference between the lowest $\chi^2$ found in each chain, same data, same code. With $\mu_0$ free (flat prior $[-0.5,+0.2]$) the posterior reaches the prior's upper edge in every combination, so the median and the lower end of the $90\,\%$ interval are given; the prediction $-0.135$ lies inside every interval (Chapter~\ref{ch:latetime}).

\observed\ What the chains show:
\begin{enumerate}
 \item With the coupling fixed by theory and no parameter added, IAM is consistent with the CMB in every data combination.
       $\Delta\chi^2$ between $+0.5$ and $+1.7$ is far below any threshold of preference either way.
 \item $\sigma_8$ falls by $0.0128$--$0.0129$ in every combination, and the standard parameters do not move.
       The direction is that of the weak-lensing surveys (KiDS-1000, DES Y3, HSC Y3~\cite{Asgari2021,DESY3,HSCY3}).
 \item When $\mu_0$ is left free, Planck alone cannot choose: the posteriors are consistent with zero ($0.05$--$0.31\sigma$)
       and with the prediction $\mu_0=-0.135$ ($0.8$--$1.4\sigma$).
 \item Placing the term in the background instead (Level 2b) drives $H_0$ to $61.45$ and $61.52$, far from every measurement.
       The term acts on matter perturbations, not on the expansion.
 \item The Level 2 code, which writes the mechanism itself into CAMB, reproduces the Level 1 result ($\Delta\chi^2=+0.54$, $\sigma_8=0.800$).
\end{enumerate}

\section{Two Hubble rates}

The photon-sector expansion rate is the one the CMB measures, $H_0=67.16$ (Level 2). For matter, the information term adds to the
effective expansion felt by bound structure; with $\beta_m=0.15765$ the matter-sector rate is $H_0^{\mathrm{matter}}=H_0\sqrt{1+\beta_m}=67.16\times1.0759=72.26\ \mathrm{km\,s^{-1}\,Mpc^{-1}}$ \derived.
IAM places the type Ia supernova distance ladder, calibrated on matter-dominated astrophysics, on the matter ruler.
\openprob\ The ladder measures distances with light. The chapter on open problems states what must be shown: how a distance built from photons
inherits the matter-sector rate, and which observable reads $H_0^{\mathrm{matter}}$ without passing through a null-geodesic distance.

\section{The decisive test}

\prediction\ Euclid~\cite{Euclid2025} is expected to constrain $\mu_0$ to $\sigma(\mu_0)\approx0.04$. the informational term predicts $\mu_0=-0.135$ with $\Sigma=1$:
a $3.4\sigma$ departure from general relativity in matter, with no change in lensing by light.
The prediction contains no parameter adjusted to the data that will test it.
